% MA 214 Discussion 9 -- covers Lecture 13.
%
% Student handout. Page 1 is the concept review; the question set runs from page 2.
% Compile from inside this folder:  pdflatex Discussion9.tex

\documentclass[11pt]{article}
\usepackage[margin=1in]{geometry}
\usepackage{parskip}
\usepackage{fancyhdr}
\usepackage{array}
\usepackage{booktabs}
\usepackage{enumitem}
\usepackage{amsmath}
\usepackage{tabularx}
\usepackage{listings}

\pagestyle{fancy}
\fancyhf{}
\cfoot{\thepage}
\rhead{Discussion 9}
\lhead{MA 214: Applied Statistics}
\renewcommand{\headrulewidth}{.4mm}

\newcommand{\blankline}[1]{\underline{\hspace{#1}}}
\newcolumntype{L}{>{\raggedright\arraybackslash}X}

\lstset{
  basicstyle=\ttfamily\small,
  columns=fullflexible,
  frame=single,
  framerule=0.4pt,
  breaklines=true,
  showstringspaces=false,
  aboveskip=8pt,
  belowskip=8pt
}

% Section header: bold title over a hairline rule.
\newcommand{\parttitle}[1]{%
  \par\medskip\noindent
  {\large\bfseries #1}\par
  \vspace{-.5em}\noindent\rule{\linewidth}{0.4pt}\par\smallskip}

% Writing space.
\newcommand{\answerspace}[1]{\vspace{#1}}

% Flush-left question list: the label runs in at the margin and the body wraps
% back to the margin, so nothing is pushed far to the right.
\newlist{questions}{enumerate}{1}
\setlist[questions]{label=\textbf{Question \arabic*.}, wide=0pt, itemsep=1.4em, topsep=.8em}

\newlist{parts}{enumerate}{1}
\setlist[parts]{label=\textbf{(\Alph*)}, leftmargin=1.6em, labelsep=.45em,
  itemsep=.7em, topsep=.5em}

\begin{document}

\noindent\textbf{Name:}\blankline{2.3in}\hfill\textbf{BUID:}\blankline{1.6in}

\begin{center}
  {\Large\bfseries Discussion 9: The Slope Is Just Another Statistic}
\end{center}

\noindent\textbf{Goals:} \textbf{(1)} read a slope, its standard error, and its interval off
\texttt{lm()} output, \textbf{(2)} say what gets shuffled to test a slope, \textbf{(3)} test a slope
against a value other than zero, \textbf{(4)} diagnose which LINE condition has failed,
\textbf{(5)} apply a transformation and read the new slope.

\parttitle{Part 1: Concept Review}

{\small

\noindent Everything in Part 2 can be answered from this page. Keep it in front of you.

\begin{enumerate}[label=\textbf{\arabic*.}, wide=0pt, itemsep=.4em, topsep=.5em]

\item \textbf{The slope has a sampling distribution.} $\hat\beta_1$ is computed from data, so a new
sample gives a different value. Everything from Module 2 applies to it unchanged: it has a sampling
distribution, a standard error, a null distribution, and a confidence interval.

\item \textbf{Two roads, now for the slope.}

\begin{center}
\renewcommand{\arraystretch}{1.15}
\begin{tabular}{lll}
\toprule
 & \textbf{Road A: simulate} & \textbf{Road B: formula} \\
\midrule
Test $H_0: \beta_1 = 0$ & shuffle the \emph{pairing} of $y$ to $x$ & $t = \hat\beta_1/\text{SE}$, df $= n-2$ \\
Interval & bootstrap the $(x,y)$ \emph{pairs} & $\hat\beta_1 \pm t^{\ast}\,\text{SE}$ \\
\bottomrule
\end{tabular}
\end{center}

Shuffling $y$ against a fixed $x$ destroys any relationship while keeping both marginal
distributions, which is exactly what $\beta_1 = 0$ claims. For the interval you must resample whole
\emph{pairs}; resampling $x$ and $y$ separately would test the null instead of estimating.

\item \textbf{Reading the output.} From an \texttt{lm()} coefficient table,
\[
t = \frac{\hat\beta_1}{\text{SE}},
\qquad
\text{95\% CI} = \hat\beta_1 \pm t^{\ast}\,\text{SE},
\qquad \text{df} = n-2 .
\]
The reported $p$-value always tests $H_0: \beta_1 = 0$, which is often not the question you care
about. The intercept is the fitted value at $x = 0$; if $x = 0$ lies outside the data, do not
interpret it.

\item \textbf{Testing against a value other than zero.} To test $H_0: \beta_1 = c$, build the
statistic yourself:
\[
t = \frac{\hat\beta_1 - c}{\text{SE}}, \qquad \text{df} = n-2 .
\]
Equivalently, ask whether $c$ lies inside the confidence interval. A slope can be
overwhelmingly different from $0$ and perfectly consistent with $c = 1.5$ at the same time.

\item \textbf{LINE.} Road B's formula needs: \textbf{L}inear mean response, \textbf{I}ndependent
observations, \textbf{N}ormal residuals, \textbf{E}qual variance. Check L and E on a plot of
residuals against fitted values, and N on a Q--Q plot. A \emph{funnel} in the residual plot means E
has failed; a \emph{curve} means L has failed.

\item \textbf{When a condition fails.} Comparing the two roads is \emph{not} a reliable way to
detect it: a $t$ interval for a slope is fairly robust, so the two roads can agree closely while a
condition is plainly broken. Diagnose from the residuals, not from agreement. The usual remedy for
a funnel is to transform $y$; if the relationship is multiplicative, take logs of both. A
\textbf{log--log} slope is an elasticity: a $1\%$ rise in $x$ goes with a $\hat\beta_1\%$ rise in
$y$, and more precisely $y$ multiplies by $1.10^{\hat\beta_1}$ when $x$ rises $10\%$.

\item \textbf{Four traps.} Reading the printed $p$-value as a test of the slope you actually care
about. Interpreting an intercept outside the data. Bootstrapping $x$ and $y$ separately when you
wanted an interval. Concluding the conditions are fine because Road A and Road B agree.

\end{enumerate}

}

\newpage

\parttitle{Part 2: Question Set}

\noindent\textbf{Scenario A.} For $n = 28$ neighbourhood branches, annual visits (thousands) are
regressed on weekly open hours. Open hours run from $22$ to $67$.

\begin{center}
\renewcommand{\arraystretch}{1.2}
\begin{tabular}{lccc}
\toprule
 & \textbf{Estimate} & \textbf{Std. Error} & \textbf{$p$-value} \\
\midrule
(Intercept) & $19.695$ & 5.383 & 0.0011 \\
\texttt{hours} & $1.3548$ & 0.1149 & $<0.001$ \\
\bottomrule
\end{tabular}
\qquad $R^2 = 0.843$, \quad $t^{\ast}_{26} = 2.0555$
\end{center}

\begin{questions}

%%%%%%%%%%%%%%%%%%%%%%%%%%%%%%%%%%%%%%
\item \textbf{Reading the slope.}

\begin{parts}
\item Verify the $t$ statistic for \texttt{hours} from the estimate and standard error, and state
its degrees of freedom.
\answerspace{3.5em}

\item Compute the 95\% confidence interval for the slope.
\answerspace{4em}

\item Interpret both the slope and its interval in context, with units.
\answerspace{4em}

\item The intercept has $p = 0.0011$, so it is ``significant.'' Explain why you would still refuse
to interpret it.
\answerspace{3.5em}
\end{parts}

%%%%%%%%%%%%%%%%%%%%%%%%%%%%%%%%%%%%%%
\item \textbf{Two roads for the slope.} Road A shuffled the pairing $1000$ times; \textbf{none} of
the $1000$ shuffled slopes was as far from $0$ as $1.3548$. Bootstrapping the pairs $2000$ times
gave a percentile interval of $[\,1.168,\; 1.549\,]$.

\begin{parts}
\item Exactly what is shuffled, and why does that operation encode $H_0: \beta_1 = 0$? Say what
stays fixed.
\answerspace{4em}

\item Report Road A's $p$-value. It is not $0$; explain what you write instead.
\answerspace{3em}

\item Compare the bootstrap interval with your $t$ interval from Question 1(B). Do they agree?
\answerspace{3.5em}

\item For the interval you must resample whole $(x,y)$ pairs. What would you be computing instead
if you resampled \texttt{hours} and \texttt{visits} separately?
\answerspace{3.5em}
\end{parts}

%%%%%%%%%%%%%%%%%%%%%%%%%%%%%%%%%%%%%%
\item \textbf{A different null value.} The director's budget model assumes each extra weekly open
hour brings exactly $1.5$ thousand annual visits. She wants that assumption tested.

\begin{parts}
\item Build the appropriate $t$ statistic and compute it.
\answerspace{4em}

\item With df $= 26$, this gives $p = 0.217$. State the conclusion about the director's assumption.
\answerspace{3.5em}

\item The same fit rejects $H_0: \beta_1 = 0$ with $p < 0.001$ and fails to reject
$H_0: \beta_1 = 1.5$ with $p = 0.217$. Explain to the director how both can be true at once.
\answerspace{4.5em}

\item Answer (B) again without any new arithmetic, using only your interval from Question 1(B).
\answerspace{3em}
\end{parts}

%%%%%%%%%%%%%%%%%%%%%%%%%%%%%%%%%%%%%%
\item \textbf{Scenario B: a condition fails.} A second regression uses \texttt{pop}, the population
served in thousands (range $9.6$ to $85.5$), to predict visits (range $10.5$ to $310.2$).

\begin{center}
\renewcommand{\arraystretch}{1.2}
\begin{tabular}{lcc}
\toprule
 & \textbf{Estimate} & \textbf{Std. Error} \\
\midrule
(Intercept) & $1.073$ & 22.629 \\
\texttt{pop} & $2.4939$ & 0.5661 \\
\bottomrule
\end{tabular}
\qquad $R^2 = 0.427$
\end{center}

\noindent Diagnostics: the residual standard deviation is $36.2$ among branches below the median
population and $77.9$ among those above it. The formula interval for the slope is
$[\,1.330,\; 3.657\,]$ and the bootstrap interval is $[\,1.131,\; 3.692\,]$.

\begin{parts}
\item Which letter of LINE has failed, and which diagnostic number told you?
\answerspace{3.5em}

\item Sketch, in words, what the residuals-against-fitted plot looks like here.
\answerspace{3em}

\item The two intervals have widths $2.33$ and $2.56$, within about $10\%$ of each other. A
classmate concludes the conditions must be fine. Explain why that reasoning fails, and say what you
should have looked at instead.
\answerspace{4.5em}

\item Even setting the interval aside, why is ``each extra thousand people brings $2.49$ thousand
visits'' a poor description of this relationship?
\answerspace{4em}
\end{parts}

%%%%%%%%%%%%%%%%%%%%%%%%%%%%%%%%%%%%%%
\item \textbf{The remedy, and find the bugs.} Refitting with both variables logged gives a
\texttt{log(pop)} slope of $0.9567$ with standard error $0.1942$, an interval of
$[\,0.558,\; 1.356\,]$, and residual standard deviations of $0.595$ and $0.688$ in the two halves.
$R^2$ rises to $0.483$.

\begin{parts}
\item Compare the two residual standard deviations with those in Question 4. Did the transformation
do its job?
\answerspace{3em}

\item Interpret the log--log slope for a $10\%$ increase in population served. You may use
$1.10^{0.9567} = 1.0955$.
\answerspace{3.5em}

\item A classmate wrote this randomization test for Scenario A's slope. It runs, prints a plausible
number, and is wrong in \textbf{three} separate ways.

\begin{lstlisting}[language=R]
obs <- coef(lm(visits ~ hours))[2]

perm <- replicate(1000, {
  vs <- sample(visits, replace = TRUE)
  coef(lm(visits ~ hours))[2]
})

mean(perm > obs)
\end{lstlisting}

Bug 1 is inside \texttt{coef(lm(...))}. What does every one of the $1000$ values equal, and what
does the last line then print?
\answerspace{4em}

\item Bug 2 is \texttt{replace = TRUE} and bug 3 is the last line. Name both, then rewrite the
whole block correctly.
\answerspace{5.5em}
\end{parts}

\end{questions}

\end{document}
